% ------------------------------------------------------------------
%  Capitolo 5 — Il cervello che impara
%  Parte II
%  Ricerca di riferimento: ricerca 01 §3
%  Voci bibliografiche principali: cowan2001, baddeley1974, craik1972, tulving1973, kandel2001, nader2000
%  Epigrafe: ✔ verificata sul testo (vedi ricerca 03 e registro, ricerca 00)
%  Scheda del capitolo: ricerca/06_indice-ragionato.md, capitolo 5
% ------------------------------------------------------------------
\chapter{Il cervello che impara}
\label{cap:cervello}

\epigrafe[Grande è questa potenza della memoria, troppo grande, Dio mio.]{Magna ista vis est memoriae, magna nimis, deus meus.}{Agostino, Confessiones X, 8, 15}

\iniziale{Q}{uesto capitolo} \segnaposto{apertura: da scrivere — vedi la scheda nell'indice ragionato}.

\section{L'attenzione, porta d'ingresso}

\segnaposto{da scrivere}

\section{La memoria di lavoro, piccola e preziosa}

\segnaposto{da scrivere}

\section{La memoria a lungo termine}

\segnaposto{da scrivere}

\section{Codificare, consolidare, ritrovare}

\segnaposto{da scrivere}

\section{Dimenticare non è un difetto}

\segnaposto{da scrivere}

\begin{daricordare}
\segnaposto{il principio del capitolo in due righe}
\end{daricordare}

\begin{domande}
  \item \segnaposto{domanda di recupero su questo capitolo}
  \item \segnaposto{domanda di applicazione a una materia che studi}
  \item \segnaposto{domanda di ripresa su un capitolo precedente}
\end{domande}

\begin{sintesi}
  \item \segnaposto{punto di sintesi}
\end{sintesi}

\finecapitolo
